\backhead{FURTHER EXERCISES}{khred}
\begin{exlist}
\item A crop with $T_b=5$ $^\circ$C needs 1500 $^\circ$C\,d to maturity. With mean temperatures of 14 $^\circ$C in May, 18 in June, 21 in July and 20 in August (31, 30, 31, 31 days), on which month does it mature?\exans{in August (1165 $^\circ$C\,d by the end of July)}
\item Find the leaf area that intercepts 90 per cent of the light with $k=0.5$ and with $k=0.8$.\exans{4.61; 2.88}
\item A maize crop with RUE 3.8 g\,MJ$^{-1}$ intercepts 85 per cent of 11 MJ\,m$^{-2}$\,d$^{-1}$ of PAR for 70 days. Find the dry matter produced and the grain yield with a harvest index of 0.5.\exans{2487 g\,m$^{-2}$; 12.4 Mg\,ha$^{-1}$}
\item A crop of 1200 g\,m$^{-2}$ has $c_m=0.012$ g\,g$^{-1}$\,d$^{-1}$ at 25 $^\circ$C and $Q_{10}=2$. Find the maintenance respiration at 18 $^\circ$C.\exans{8.9 g\,m$^{-2}$\,d$^{-1}$}
\item A soybean crop ($K_y=0.85$) receives 70 per cent of its seasonal water need. With $Y_m=3.5$ Mg\,ha$^{-1}$, find $Y_a$.\exans{2.61 Mg\,ha$^{-1}$}
\end{exlist}
\backhead{ESSAY QUESTIONS}{khred}
{\small\itshape Answer each question in about one page, with equations, a sketch and at least one number from this chapter.}\par\vspace{3pt}
\begin{exlist}
\item Describe the chain of stores of the crop in Daisy, from the light to the harvested organ.
\item Explain thermal time and why it is a better clock for a crop than calendar days.
\item Discuss the light-response curve of a leaf and why a canopy is more efficient than its leaves.
\item Explain maintenance and growth respiration and the effect of warming on the carbon budget of a crop.
\item Discuss how water stress enters the crop model and how it is expressed in yield.
\item Explain the management column of the process map and how sowing, fertilisation, irrigation and harvest enter a simulation.
\end{exlist}
\dykband{EXERCISES IN SWAFL}
{\small Open SWAFL (desktop or online) with the benchmark project. The chapter of the \href{https://github.com/halilunaziru73-creator/Coupled-Water-Fluxes-Pore-Scale-Dynamics-Precision-Irrigation-Mulching-Feedback-Loops/raw/main/Software/v1.1.0/SWAFL_Desktop/docs/manual_v1.1.0/SWAFL_Manual_v1.1.0.pdf}{SWAFL manual} that explains each task is given in brackets; record your values in a table and compare them with what is expected (\emph{Check}).}\par\vspace{3pt}
\begin{exlist}
\item {\color{khblue}[Manual, Chapter 25]} Click every box of the crop part of the Daisy process map and list the stores and the flows of carbon.\par\noindent{\small\itshape Check: leaves, stems, storage organ and roots.}
\item {\color{khblue}[Manual, Chapter 30]} Open the dry matter of the organs of a run and read the harvest index of three years.\par\noindent{\small\itshape Check: the harvest index varies with the season.}
\item {\color{khblue}[Manual, Chapter 21]} Change the crop of the project and compare the crop water requirement.\par\noindent{\small\itshape Check: each crop has its own $K_c$ curve.}
\item {\color{khblue}[Manual, Chapter 28]} Design a Daisy experiment that varies the sowing date over four levels and run it.\par\noindent{\small\itshape Check: the yield changes with the sowing date.}
\end{exlist}
\backhead{OBJECTIVE TYPE QUESTIONS}{khred}
\begin{exlist}
\item Thermal time counts degrees\opt{below $T_b$}{above $T_b$}{above 30 $^\circ$C}{at night only}
\item A canopy with $kL=3$ intercepts about\opt{50}{75}{95}{100} per cent of the light
\item A $Q_{10}$ of 2 means the rate\opt{halves per 10 $^\circ$C}{doubles per 10 $^\circ$C}{doubles per 1 $^\circ$C}{is constant}
\item $K_y>1$ means the yield falls\opt{slower}{as fast as}{faster than}{independently of} the water use
\item The conversion efficiency of growth is about\opt{0.2}{0.5}{0.72}{1.0}
\end{exlist}
\objkey{1.\ (b),\quad 2.\ (c),\quad 3.\ (b),\quad 4.\ (c),\quad 5.\ (c)}
